\section{Conclusion}\label{sec:conclusion}

We presented, in full-wave simulation, the calibrated posterior: a
flow-matching/generative WCE-RF position posterior with empirical multi-level
coverage under position-grouped cross-validation at a four-receiver MICS/MedRadio
budget in the legal MICS/MedRadio band \SIrange{401}{406}{MHz}. It is the only
evaluated method that is simultaneously accurate, well covered from its own
samples (conservative at 50\%, near-nominal at 90/95\%), and discriminative in
error-vs-spread. Its posterior mean also preserves deterministic-backbone point
accuracy within \SI{0.014}{cm} (\SI{0.354}{cm} vs \SI{0.340}{cm}) and beats ten
classical regressors on identical position-grouped folds. Beyond aggregate
calibration, we showed two properties unique to a generative posterior: in a
controlled mirror-symmetric receiver-array ablation its \emph{modality} tracks
forward-operator invertibility, and on the main dataset its \emph{width} carries
the aleatoric uncertainty of an unknown capsule orientation. A
feature-deployability ablation places the
single-tone magnitude-only floor at ${\sim}\SI{1.5}{cm}$ and attributes the
sub-centimetre regime to phase coherence rather than bandwidth.

We make no claim to state-of-the-art point accuracy; real-phantom systems at
higher frequencies are more accurate, and we declare the simulation-vs-real and
band confounds throughout. To our knowledge, the novelty is the first
flow-matching/generative WCE-RF position posterior with empirical multi-level
coverage under position-grouped cross-validation. The next evidence gate is a
physical-phantom companion that tests whether the posterior remains calibrated on
\emph{measured} channels under the same position-grouped protocol. We also flag
a true $SO(3)$ orientation marginal and paired deterministic-backbone predictions
for a direct CFM-vs-backbone interval as concrete next steps.
